%% ============================================================
\section{Introduction}
\label{sec:intro}
%% ============================================================

\subsection{Motivation and Problem Statement}

Software development has undergone a structural transformation over the past
decade. The rise of geographically distributed teams, remote-first engineering
cultures accelerated by the COVID-19 pandemic~\cite{ford2019remote}, and the
emergence of \textit{AI pair-programming} tools such as GitHub
Copilot~\cite{chen2021evaluating} and Cursor have collectively redefined what
developers expect from their tooling.  Modern practitioners no longer regard
the Integrated Development Environment (IDE) as a purely local, single-user
artefact; they demand \emph{collaborative, cloud-native workspaces} that
synchronise edits in real time, facilitate async and synchronous communication,
and harness Large Language Models (LLMs) to accelerate routine programming
tasks.

Despite significant commercial investment, the current landscape of
collaborative coding tools exhibits three persistent gaps:

\begin{enumerate}[leftmargin=*, label=\textbf{G\arabic*.}]
  \item \textbf{Siloed Functionality.} Tools such as Visual Studio Live Share
    provide code synchronisation but require a local IDE installation and lack
    integrated project management, persistent file storage, or AI assistance.
    Conversely, AI coding assistants (Copilot, CodeWhisperer~\cite{amazon2023codewhisperer})
    operate within individual developer contexts and offer no collaborative
    workspace.

  \item \textbf{Inadequate Access-Control Granularity.} Platforms such as
    Replit and CodeSandbox offer collaborative editing but treat all
    collaborators uniformly, precluding fine-grained ownership semantics
    essential for enterprise and academic use cases where a designated
    \emph{project leader} must control team membership and project lifecycle.

  \item \textbf{Weak AI-Collaboration Coupling.} Existing tools append AI as
    an afterthought---a side-panel chat disconnected from the shared workspace.
    No widely-deployed open-source platform broadcasts AI-generated code
    artefacts in real time to \emph{all} project collaborators simultaneously,
    preserving shared context across the team.
\end{enumerate}

\subsection{Current Gap in Research}

Recent literature (2024--2026) highlights the growing maturity of LLM-assisted
programming. Wei \etal~\cite{wei2022chain} introduced chain-of-thought
prompting, enabling models to reason step-by-step through complex coding
tasks. Concurrent work by Peng \etal~\cite{peng2023impact} demonstrated
a~55.8\% acceleration in task completion for programmers using AI assistance.
Jiang \etal~\cite{jiang2024survey} surveyed 67 LLM-for-code studies and
identified \emph{context sharing} as a primary unsolved challenge: when a team
member invokes an AI assistant, the generated code lives only in their local
context and must be manually shared.

On the real-time systems side, studies of operational-transformation
(OT)~\cite{sun1998operational} and conflict-free replicated data types
(CRDTs)~\cite{shapiro2011conflict} underpin collaborative text editors such as
Google Docs~\cite{googledocs2010} and Firepad. However, the majority of
open-source implementations focus exclusively on text synchronisation, leaving
file-tree management, binary artefacts, and project-scoped permission models
as open engineering problems.

From a security perspective, Almohri \etal~\cite{almohri2023security} and
recent OWASP guidance on token management~\cite{owasp2024jwt} emphasise the
inadequacy of stateless JWT invalidation---a gap that cloud-native platforms
typically address through in-memory or distributed-cache token blacklists, an
approach that \echoray systematically implements via Redis.

Synthesising these threads, we identify a clear and actionable research
opportunity: the design and rigorous evaluation of a \emph{unified},
\emph{open-source} platform that integrates real-time code collaboration,
role-based project governance, and AI-broadcast assistance within a single
coherent architecture.

\subsection{Novel Contributions}

This paper makes the following contributions to the literature on software
engineering tools and collaborative systems:

\begin{enumerate}[leftmargin=*, label=\textbf{C\arabic*.}]
  \item \textbf{Unified Architecture.} We present the complete design of
    \echoray---a full-stack TypeScript platform coupling a Next.js~15 /
    React~19 front end with an Express~5 / Prisma / PostgreSQL back end,
    Socket.io real-time layer, Redis session management, and Google Gemini~2.5
    Flash AI integration (\cref{sec:architecture}).

  \item \textbf{AI-Broadcast Protocol.} We define and implement a novel
    \emph{AI-broadcast} pattern in which a single \texttt{@ai} invocation in
    the project chat triggers a Gemini query whose structured JSON response is
    simultaneously delivered to all online collaborators in the same Socket.io
    room (\cref{sec:ai}).

  \item \textbf{RBAC-Gated Collaboration.} We describe a role-based access
    model distinguishing \emph{leaders} from \emph{members} enforced at the
    service layer, covering collaborator addition/removal, project deletion, and
    file-tree ownership (\cref{sec:security}).

  \item \textbf{Empirical Latency Characterisation.} We provide reproducible
    measurements of WebSocket message latency, AI response latency, and
    database query performance under multi-user load, establishing baseline
    benchmarks for future research (\cref{sec:evaluation}).

  \item \textbf{Open-Source Release.} The complete source code is publicly
    available, enabling replication, extension, and community-driven
    improvement.
\end{enumerate}

\subsection{Conceptual Overview}

\Cref{fig:overview} presents a high-level conceptual diagram of \echoray,
illustrating the three principal interaction planes: (1)~the \emph{User
Interaction Plane} encompassing browser-based code editing, chat, and project
management; (2)~the \emph{Real-Time Synchronisation Plane} mediated by
Socket.io over WebSocket; and (3)~the \emph{AI Generation Plane} bridging
natural-language chat prompts to LLM-generated code artefacts broadcast back
to all collaborators.

\begin{figure}[H]
  \centering
  \begin{tikzpicture}[
    font=\small,
    node distance=0.5cm and 1.2cm,
    box/.style={
      rectangle, rounded corners=4pt,
      minimum width=2.6cm, minimum height=0.75cm,
      text centered, align=center, draw, thick
    },
    plane/.style={
      rectangle, rounded corners=6pt,
      draw, dashed, thick,
      inner sep=8pt
    },
    arr/.style={-Stealth, thick},
    bidir/.style={Stealth-Stealth, thick},
  ]

  %% ── User Interaction Plane ────────────────────────────────
  \node[box, fill=blue!15]  (editor)  {Monaco Editor};
  \node[box, fill=blue!15, right=of editor]  (chat)    {Chat Box};
  \node[box, fill=blue!15, right=of chat]    (filetree){File Tree};
  \node[box, fill=blue!15, right=of filetree](preview) {Live Preview};

  \begin{scope}[on background layer]
    \node[plane, fit=(editor)(chat)(filetree)(preview),
          label=above:\textbf{User Interaction Plane (Next.js / React)}]
      (uiplane) {};
  \end{scope}

  %% ── Real-Time Synchronisation Plane ───────────────────────
  \node[box, fill=green!20, below=1.4cm of editor]  (socket)  {Socket.io\\Server};
  \node[box, fill=green!20, right=of socket]         (redis)   {Redis\\(Token Store)};
  \node[box, fill=green!20, right=of redis]          (express) {Express 5\\REST API};
  \node[box, fill=green!20, right=of express]        (prisma)  {Prisma ORM\\(PostgreSQL)};

  \begin{scope}[on background layer]
    \node[plane, fit=(socket)(redis)(express)(prisma),
          label=above:\textbf{Real-Time Synchronisation \& API Plane (Node.js)}]
      (rtplane) {};
  \end{scope}

  %% ── AI Generation Plane ───────────────────────────────────
  \node[box, fill=orange!25, below=1.4cm of socket]  (aisvc)   {AI Service\\(ai.service.ts)};
  \node[box, fill=orange!25, right=of aisvc]          (gemini)  {Google Gemini\\2.5 Flash};
  \node[box, fill=orange!25, right=of gemini]         (jsonp)   {JSON Response\\Parser};

  \begin{scope}[on background layer]
    \node[plane, fit=(aisvc)(gemini)(jsonp),
          label=above:\textbf{AI Generation Plane (Google Generative AI SDK)}]
      (aiplane) {};
  \end{scope}

  %% ── Arrows: UI ↔ RT ──────────────────────────────────────
  \draw[bidir] (editor.south)   -- ++(0,-0.5) -| (socket.north);
  \draw[bidir] (chat.south)     -- ++(0,-0.3) -| (socket.north);
  \draw[bidir] (filetree.south) -- ++(0,-0.3) -| (express.north);
  \draw[arr]   (preview.south)  -- ++(0,-0.3) -| (prisma.north);

  %% ── Arrows: RT ↔ AI ──────────────────────────────────────
  \draw[arr]   (socket.south)   -- (aisvc.north)
               node[midway, right=2pt]{\textit{@ai trigger}};
  \draw[bidir] (aisvc.east)     -- (gemini.west);
  \draw[arr]   (gemini.east)    -- (jsonp.west);
  \draw[arr]   (jsonp.north)    -- ++(0,0.4) -| (socket.south)
               node[near start, right=2pt]{\textit{broadcast}};

  \end{tikzpicture}
  \caption{Conceptual overview of \echoray{}. Three interaction planes---User
    Interaction, Real-Time Synchronisation, and AI Generation---are connected
    by bidirectional WebSocket channels and REST API calls. An \texttt{@ai}
    trigger in the chat box propagates through the Socket.io server to the
    Google Gemini LLM; the structured response is broadcast to every
    collaborator in the project room.}
  \label{fig:overview}
\end{figure}

\subsection{Paper Organisation}

The remainder of this paper is organised as follows.
\Cref{sec:background} surveys related work in collaborative development
environments, real-time synchronisation protocols, and AI-assisted
programming.
\Cref{sec:architecture} presents the system architecture of \echoray{} in
detail, including the layered component model, database schema, and deployment
topology.
\Cref{sec:realtime} analyses the real-time collaboration engine, covering the
Socket.io room model, file-tree synchronisation protocol, and chat system.
\Cref{sec:ai} describes the AI integration pipeline from prompt ingestion to
broadcast delivery.
\Cref{sec:security} details the security subsystem, encompassing
authentication, token blacklisting, RBAC, and input validation.
\Cref{sec:implementation} discusses implementation decisions and the full
technology stack.
\Cref{sec:evaluation} reports empirical latency and throughput measurements.
\Cref{sec:discussion} reflects on lessons learned, limitations, and
implications for the field.
\Cref{sec:conclusion} concludes with a summary and roadmap for future work.
